\documentclass[UTF8,fontset=none,11pt,a4paper]{ctexart}
\usepackage[left=21mm,right=21mm,top=22mm,bottom=22mm,headheight=15pt,headsep=8mm]{geometry}
\usepackage{fontspec}
\setmainfont{texgyretermes-regular.otf}[BoldFont=texgyretermes-bold.otf,ItalicFont=texgyretermes-italic.otf]
\setCJKmainfont{Songti.ttc}[Path=/System/Library/Fonts/Supplemental/,FontIndex=6,AutoFakeBold=2]
\setCJKsansfont{Songti.ttc}[Path=/System/Library/Fonts/Supplemental/,FontIndex=6,AutoFakeBold=2]
\usepackage[table]{xcolor}
\definecolor{TableHead}{HTML}{EEF0ED}
\definecolor{Muted}{HTML}{666C68}
\usepackage{booktabs,longtable,array,fancyhdr,needspace}
\usepackage{xurl}
\usepackage[unicode,hidelinks]{hyperref}
\urlstyle{same}
\linespread{1.32}
\setlength{\parindent}{2em}
\setlength{\parskip}{3pt}
\setlength{\LTpre}{8pt}\setlength{\LTpost}{10pt}
\ctexset{section={format=\large\bfseries,beforeskip=18pt,afterskip=8pt},subsection={format=\normalsize\bfseries,beforeskip=12pt,afterskip=6pt}}
\pagestyle{fancy}\fancyhf{}
\renewcommand{\headrulewidth}{.3pt}
\renewcommand{\footrulewidth}{0pt}
\fancyhead[L]{\small\color{Muted}住房贷款付款服务及账户办理文件}
\fancyhead[R]{\small\color{Muted}12}
\fancyfoot[C]{\small\color{Muted}\thepage}
\hypersetup{pdftitle={住房贷款付款服务及账户办理文件},pdfauthor={金融决策模拟材料}}
\begin{document}
{\noindent\LARGE\bfseries 住房贷款付款服务及账户办理文件\par}
\vspace{7pt}{\noindent\small\color{Muted}中文模拟参考文件\par}\vspace{10pt}
\section*{第一部分　第三方付款服务登记及授权条款}
\subsection*{1. 扣款指示}
服务按登记的付款日历，每隔两周从指定存款账户扣取正常月供的一半。全年日历列示26次半额扣款。每次另收交易费2.50美元。扣款银行、账户号码及付款授权由登记资料确定。


\subsection*{2. 服务登记费用}
加入服务的登记费用为295美元。登记费用从计划产生的额外半额付款中收取；尚未收足前，额外付款先用于登记费用，余额才按安排转付本金。


\subsection*{3. 资金保管及月供转付}
从存款账户扣款，并不同时将该笔款项应用于房贷本金。服务暂存所收半额款项，按正常月供日期转付完整月供。贷款机构接收并按贷款账户规则应用付款。


\subsection*{4. 三次扣款月份}
某月发生三次半额扣款时，超出正常月供所需金额的半额款项为额外付款。按费用收取顺序，先支付尚未收完的登记费用，再把剩余款项用于下一次正常付款安排中的本金付款。


\subsection*{5. 贷款账户与服务账户}
正常月供义务仍由原贷款合同确定。服务负责收集和转付资金，贷款机构负责贷款账目。所收交易费用与登记费用不作为贷款本金付款，也不改变合同利率或贷款期限。


\subsection*{6. 本次授权记录}
登记的半额款为500美元，每次授权扣款总额502.50美元。贷款机构的月供到期日和额外本金入账日期，见本次账户办理确认。


\section*{第二部分　贷款机构付款资料}
\subsection*{一般付款常见问题}
我可以在线支付房贷或住房净权益账户吗？可以，从本机构或其他银行的存款账户付款。可以通过网站或手机应用设置。


自动付款有什么好处？不必每次记住或担心漏付款。通过机构设置后，因税费和保险托管变化，或利率变化而产生的付款额调整，会自动更新。


如何通过电话付款？可拨打房贷付款服务电话1-833-729-2427。


可以追加本金或托管付款吗？可以在线或通过手机应用办理一次性付款，也可以在自动付款中增加本金。完整月供应用前收到的额外本金款，会先显示为未分配资金，直至完整月供应用。


如何在线支付托管不足额？登录贷款账户，选择支付房贷，再选“仅本金／托管／费用”，输入追加托管或不足额，然后提交付款。


如何知道付款已收到？线上记录通常三日内反映付款。灵活付款中的部分款项先显示为未分配资金；凑足下一笔完整月供后，才记作贷款付款。


多付正常月供，是否用于减少本金？通常是，但账户应没有拖欠，也没有另行付款指示。


付款入账截止时间是什么？周一至周五，美国东部时间晚上7时30分前收到的付款，当天记入；线上通常三个营业日内显示。如果宽限期结束时仍未付完整应付款，可能收取逾期费。


\subsection*{房贷利息和付款如何运作？}
每月应付利息由贷款利率及剩余本金决定。我们采用上个月的本金余额计算当期应付利息。付款的其余部分按需要用于本金及托管。


\subsection*{付款日期是否影响当期应付利息？}
灵活付款中的月供按合同到期日应用。只要在宽限期内支付，实际付款日期不会改变这一期应付利息。额外付款应用于本金后，之后期间使用的本金余额相应改变。


\subsection*{付款日期是否重要？}
只要在宽限期内完成付款，不会因付款日期收取逾期费。


\subsection*{灵活付款常见问题}
每月两次与每两周一次，有什么不同？两者均帮助安排预算。每月两次由你选定两个相隔至少十日的日期，全年24次半额付款；每两周一次选定每隔一周的同一星期日，全年通常26次半额付款，两个月分别有三次。


每两周付款的两次额外半额付款会用于本金。服务机构按上月本金余额计算利息，并加上按期偿还贷款所需的本金。月供按到期日应用；在宽限期内支付，付款日不改变该期应付利息。


\subsection*{如何处理每月两次和每两周付款？}
我们保留不足支付一笔完整月供的款项，包括通过上述两种频率收到的款项，并在你的账户中列作“未分配资金”。


收到足以支付完整月供的金额后，我们才将其应用于贷款付款。


灵活付款会更快结清贷款吗？每两周付款通常全年26次半额付款，而不是12次完整付款。日历及所选星期日决定一年中约两个月有三次付款；额外半額付款用于本金，余额减少较快。


每月两次只是把月供拆成两笔，在所选日期扣款。除非另外追加本金，否则不会因此加快结清。


可以选择哪些扣款日？每月两次的日期至少间隔十日、不超过十五日；每两周付款必须是隔周同一个星期日。付款日历用于设置。


\subsection*{如果已经用其他方式付款，仍会自动扣款吗？}
住房抵押贷款：即使你已经通过分行、网上、邮寄或电话付款，我们仍会按自动付款设置再扣款，并将该自动付款应用于下一笔应付款。


住房净权益贷款及信用额度：如果你已通过上述渠道付款，我们不会再自动扣取另一笔付款。


设置灵活付款是否收费？不收设置费用。


已经预付月供的账户可以加入吗？可以，但不能已经预付超过两个月。


灵活自动付款能加入额外本金吗？可以；完整月供应用之前，追加本金款也显示为未分配资金。


所有住房贷款都能使用这些选项吗？不是。住房净权益贷款及信用额度、选项ARM、只付利息贷款、每日单利贷款、建设贷款、房贷优惠利率贷款及按年或季度付款的贷款，不能设置每月两次或每两周付款，仍可设置每月自动付款。


\subsection*{灵活付款产品页面}
每月一次：可用于房贷及住房净权益贷款。每月在所选日期扣款，在到期日应用付款，有助于遵守预算。


每月两次：只用于房贷。加入首月，需支付正常月供及两次半额付款；月供拆为两次，你选择日期。半额付款暂存为未分配资金，收到完整月供后才应用。


每两周一次：只用于房贷。加入首月同样需支付正常月供和两次半额付款。通常一年有两个月分别产生额外半额付款，两笔均用于本金，相当于全年多一笔完整付款。


\subsection*{房贷付款示例}
原产品页面示例使用250,000美元、4.5\%固定利率、三十年贷款，另假设每月税费与保险475美元。


每月一次支付1,742美元：三十年结清，利息206,016美元。每月两次各支付871美元：仍在三十年结清，利息206,016美元。每两周一次支付871美元：约24.5年结清，利息162,094美元；比前两种方式少约5.5年，利息少超过43,000美元。


上述金额为产品页面的另外一个示例，所用本金、利率、期数及税费保险均已单独列明。


所有住房信贷产品均须符合信用与物业审批要求。利率、项目条款和条件可能改变；并非所有产品在所有州或所有金额范围都可取得，其他限制也可能适用。


\section*{本次账户资料及办理记录}
\subsection*{A：第三方还款服务}
\textbf{双周正常扣款部分：} 每次500美元


\textbf{双周交易费：} 每次2.50美元，另外扣取


\textbf{加入费：} 295美元，从额外扣款中收取


\textbf{每年双周扣款次数：} 26次


\textbf{扣款后的资金位置：} 第三方保管账户


\textbf{正常月供入账：} 按原合同每月到期日向贷款机构支付1,000美元


\textbf{三次扣款月份的额外500美元：} 先清偿尚未收完的加入费，剩余金额在当月正常月供入账后冲减本金


\textbf{交易费与本金：} 交易费支付给服务商，不计入还贷本金


\textbf{年末的补充操作：} 本年度预算已按扣款安排使用


\textbf{额外本金付款收费：} 无提前还款罚金


\textbf{服务管理：} 向第三方登记扣款账户并查看服务扣款记录


\subsection*{B：贷款机构直接扣款}
\textbf{双周正常扣款部分：} 每次500美元


\textbf{双周交易费：} 无该交易费


\textbf{加入费：} 无加入费


\textbf{每年双周扣款次数：} 26次


\textbf{扣款后的资金位置：} 贷款机构的未分配付款记录


\textbf{正常月供入账：} 按相同到期日完成1,000美元正常付款


\textbf{三次扣款月份的额外500美元：} 在同一时点将额外500美元冲减本金


\textbf{交易费与本金：} 本次预算中保留的对应资金留在原银行账户


\textbf{年末的补充操作：} 第十二个月正常付款及额外付款入账后，将保留的对应资金另作一次本金付款


\textbf{额外本金付款收费：} 同样无提前还款罚金；本次年末在线本金付款免费


\textbf{服务管理：} 在贷款机构账户设置扣款，并核对贷款入账记录


本次比较从付款设置完成后的正常月份开始，覆盖接下来的连续十二个月。参考文件中首次设置的过渡付款发生在本窗口开始之前；该过渡步骤不改变这里已确认的120,000美元起始本金、26次扣款日历及13,065美元全年预算。


\subsection*{办理确认事项}
\textbf{加入费的收取。} A不在开通时另外扣295美元。第六个月出现额外的500美元时，服务商先从中扣除全部加入费，再把余款作为额外本金付款。第十二个月的额外500美元按表中安排处理。


\textbf{扣款与入账。} 两种方案均不把每次500美元直接视作一笔完整月供。凑足正常月供后，按合同到期日办理。本次两种方案的正常付款日、额外本金入账日及到账条件均已确认，不发生迟付。


\textbf{全年资金安排。} A每次实际从银行账户扣取500美元与2.50美元交易费。B每次自动扣500美元，其余对应预算继续留在银行账户，在年末按上述安排支付给贷款机构。两种方案均使用全年13,065美元，不在预算外补缴服务费用。


\textbf{贷款计息。} 房贷采用相同的正固定利率，月利息按上个月未偿本金计算；正常月供金额不因本次额外还款而重新降低。额外本金冲减会改变后续计息本金，贷款不会在本年度内结清。


\textbf{本次办理范围。} 两种方案的自动扣款均能正常执行；不设置忘记付款、扣款失败或新增银行手续费。服务可按流程变更，但本次选择覆盖完整年度，期间不取消、不改扣款频率。

\end{document}
